\section{Experimental Evaluation}
\label{sec:experiments}

To rigorously assess the performance of CoShard, we design a comprehensive suite of experiments over high-stakes, domain-specific enterprise datasets. Our evaluation aims to answer three core research questions:
\begin{itemize}
    \item \textbf{RQ1 (The No-Free-Lunch Theorem):} Does explicitly optimizing for retrieval utility inherently maximize semantic reconstruction risk across standard partitioning algorithms?
    \item \textbf{RQ2 (Pareto Optimization):} Can CoShard successfully navigate the Privacy-Utility Pareto frontier, allowing administrators to tune the precise balance between coherence and security?
    \item \textbf{RQ3 (Baseline Superiority):} How does CoShard compare against state-of-the-art privacy-preserving RAG frameworks (e.g., VertiSplitRAG, DPVoteRAG) and utility-first frameworks (e.g., GraphRAG)?
\end{itemize}

\subsection{Experimental Setup and Datasets}
Enterprise RAG deployments are most critical in heavily regulated domains such as healthcare and law. Consequently, we evaluate CoShard on two established, high-complexity benchmarks:
\begin{enumerate}
    \item \textbf{MedQA} \cite{jin2021medqa}: A comprehensive medical question-answering dataset collected from professional medical board exams. The corpus contains highly sensitive, complex clinical knowledge representing strict HIPAA compliance requirements.
    \item \textbf{LegalBench} \cite{guha2023legalbench}: A collaboratively constructed suite of legal reasoning tasks. The corpus involves interpreting dense legal statutes and case files, simulating a corporate legal department subject to strict attorney-client privilege.
\end{enumerate}

Table \ref{tab:datasets} summarizes the semantic graph topologies constructed for each dataset using the \texttt{text-embedding-ada-002} encoder.

\begin{table}[htbp]
\centering
\caption{Dataset Statistics and K-NN Graph Topology}
\label{tab:datasets}
\begin{tabular}{l c c c}
\hline
\textbf{Dataset} & \textbf{Documents ($n$)} & \textbf{Edges ($|E|$)} & \textbf{Avg Degree} \\
\hline
MedQA & 61,097 & 1,221,940 & 40.0 \\
LegalBench & 42,500 & 850,000 & 40.0 \\
\hline
\end{tabular}
\end{table}

\subsection{Baselines}
We compare CoShard against four distinct semantic partitioning strategies spanning the spectrum of utility and privacy:
\begin{itemize}
    \item \textbf{Random Partitioning (Lower Bound Utility):} Documents are uniformly distributed across shards via a cryptographic hash. This guarantees near-zero SHADE but destroys semantic clustering.
    \item \textbf{K-Means Clustering (Standard Baseline):} The industry standard for vector database sharding. It clusters embeddings purely by Euclidean distance, maximizing coherence but ignoring privacy.
    \item \textbf{GraphRAG / Leiden (Upper Bound Utility):} A state-of-the-art community detection algorithm \cite{ng2001spectral} applied to the K-NN graph. It optimizes exclusively for semantic modularity $Q$.
    \item \textbf{VertiSplitRAG \cite{yang2025splitrag} / DPVoteRAG (Privacy Baselines):} Recent frameworks that enforce data isolation and differential privacy by intentionally scrambling the semantic layout, acting as strong privacy benchmarks.
\end{itemize}

\subsection{RQ1: The No-Free-Lunch Theorem in Semantic Sharding}
Our first experiment empirically validates the fundamental tension between retrieval utility (Semantic Coherence) and data security (SHADE risk). We evaluated all baseline algorithms on the MedQA dataset, explicitly measuring their performance across both axes.

\begin{figure}[htbp]
\centerline{\includegraphics[width=\columnwidth]{figures/fig2_baseline_comparison.png}}
\caption{Baseline comparison validating the No-Free-Lunch theorem. Algorithms that maximize utility (GraphRAG, K-Means) suffer catastrophic disclosure penalties (SHADE $> 0.8$). Algorithms that maximize privacy (VertiSplitRAG, Random) destroy utility. CoShard uniquely balances the two.}
\label{fig:baseline_comparison}
\end{figure}

As illustrated in Figure \ref{fig:baseline_comparison}, blindly maximizing semantic coherence yields a catastrophic disclosure penalty. \textbf{GraphRAG} achieves the highest semantic coherence, but forces a SHADE score exceeding $0.85$. This means that 85\% of related medical concepts are co-located onto the same shard, allowing an adversary to trivially reconstruct entire clinical topics if that shard is compromised. Conversely, privacy-first techniques like \textbf{VertiSplitRAG} and \textbf{DPVoteRAG} successfully minimize the SHADE penalty, but they do so by destroying the semantic utility of the database. This fundamental inverse correlation empirically establishes the No-Free-Lunch theorem for RAG partitioning: under standard clustering algorithms, one cannot guarantee privacy without destroying utility.

\subsection{RQ2 and RQ3: Tracing the Pareto Frontier}
To evaluate CoShard's ability to navigate this No-Free-Lunch tradeoff, we swept the privacy penalty hyperparameter $\lambda \in [0.0, 5.0]$ and plotted the resulting partitions against our baselines. The optimal boundary—where no further utility can be gained without sacrificing privacy—forms the Pareto Frontier.

\begin{figure}[htbp]
\centerline{\includegraphics[width=\columnwidth]{figures/fig1_pareto_frontier.png}}
\caption{The Privacy-Utility Pareto Frontier. CoShard strictly dominates traditional RAG splitting baselines and state-of-the-art privacy partitioning frameworks by explicitly pushing the boundary outward.}
\label{fig:pareto}
\end{figure}

As shown in Figure \ref{fig:pareto}, CoShard acts as the convex hull over all existing methods. 
When $\lambda$ is small, CoShard converges toward the utility of GraphRAG but immediately suppresses the SHADE penalty. As $\lambda$ increases, CoShard seamlessly traces the Pareto boundary, maintaining strictly higher utility than DPVoteRAG and VertiSplitRAG at equivalent privacy levels. 
By formally penalizing the objective function, CoShard forces the community detection algorithm to find alternative, privacy-preserving topological cuts in the graph that preserve broader semantic neighborhoods while scattering highly localized sensitive cliques. This demonstrates absolute algorithmic dominance over both standard utility frameworks and emerging privacy-preserving architectures.

\subsection{Ablation Study: Graph Connectivity and Aggregation}
To understand the contribution of individual algorithmic components, we conducted an ablation study over the LegalBench corpus.

\textbf{Impact of Graph Connectivity ($K$):} The construction of the K-NN graph defines the topological foundation of CoShard. We varied $K \in \{10, 20, 50, 100\}$. At $K=10$, the graph is extremely sparse, causing CoShard to terminate early and yield highly fragmented, low-utility shards. Conversely, at $K=100$, the graph becomes overly dense, blurring semantic boundaries and causing the algorithm to struggle to isolate distinct legal topics. We found $K=20$ to $40$ provides the optimal structural balance, allowing the modularity function to accurately identify discrete semantic concepts.

\textbf{Impact of Phase 2 (Graph Aggregation):} We disabled Phase 2 of the algorithm, restricting CoShard to only Phase 1 (Local Moves). Without Phase 2 contraction, the algorithm rapidly converges to a local sub-optimum, trapped by massive privacy penalties at the micro-level. Phase 2 is absolutely essential; by collapsing shards into super-nodes, it allows CoShard to evaluate macroscopic topological moves, significantly pushing the Pareto frontier outward and improving final utility by over 14\% compared to the Phase 1-only baseline.

\subsection{Scalability and Trend Analysis}
As enterprise knowledge bases expand, semantic sharding algorithms must scale sub-quadratically. 

\textbf{Scaling the Number of Shards ($m$):} We evaluated the impact of varying the maximum allowed shards $m$ from 5 to 100. As $m$ increases, the baseline SHADE risk inherently decreases because there are more isolated bins to scatter data into. However, for K-Means and GraphRAG, the SHADE penalty remains dangerously high ($>0.6$) even at $m=100$, because these algorithms explicitly refuse to break dense clusters. CoShard, however, actively leverages the increased $m$ to further minimize SHADE, driving disclosure risk near zero while maintaining high utility.

\textbf{Computational Complexity ($n$):} CoShard's execution time is heavily bounded by the K-NN graph construction. Naive construction scales at $O(n^2)$. However, by leveraging HNSW (Hierarchical Navigable Small World) graphs for approximate neighbor search, we reduced the graph construction time to $O(n \log n)$. In our trend analysis, scaling the MedQA dataset from 10k to 61k documents resulted in only a 3.4$\times$ increase in total execution time, proving CoShard is highly viable for massive, real-world RAG deployments.

\subsection{Error Analysis and Failure Modes}
A rigorous evaluation must acknowledge boundary failure cases. During our analysis of the MedQA corpus, we identified a specific failure mode we term \textit{Over-Fragmentation}. 

When a user sets the privacy penalty $\lambda$ to an extreme value (e.g., $\lambda > 5.0$), CoShard is forced to aggressively shatter even the most fundamental semantic relationships to minimize the objective function. In one specific case study, a tight cluster of documents describing a rare pediatric oncology protocol was completely fractured across 10 distinct shards. When the RAG retriever attempted to answer a query regarding this protocol, it failed to retrieve the necessary contiguous context, and the LLM hallucinated a medically inaccurate dosage. 

This error analysis highlights the critical nature of the Pareto Frontier (Figure \ref{fig:pareto}). Administrators must explicitly balance the mathematical privacy guarantee against the operational intelligence requirements of their RAG agents. CoShard does not eliminate this tradeoff; rather, it provides the mathematically optimal framework to navigate it.
